\chapter{글로벌 AI 거버넌스 프레임워크 비교 분석}

AI 거버넌스는 더 이상 개별 기업이나 국가 차원의 문제가 아니다. AI 시스템이 국경을 초월하여 개발·배포·운영되는 현실에서, 글로벌 수준의 거버넌스 프레임워크에 대한 이해는 한국 기업의 필수 역량이 되었다.\cite{fjeld2020principled} 본 장에서는 국제기구와 표준화 기관이 제시하는 주요 프레임워크를 분석하고, 이들이 한국 기업의 AI 거버넌스 체계 수립에 주는 시사점을 도출한다. 나아가 주요국의 규제 동향을 심층 분석하고, 다국적 기업이 활용할 수 있는 통합 규제 대응 전략과 실무 도구를 제시한다.

\section{국제기구 및 표준화 기관의 프레임워크}
국제기구와 표준화 기관들이 제시하는 AI 거버넌스 프레임워크는 크게 두 가지 유형으로 구분된다. 첫째, OECD와 UNESCO가 대표하는 원칙 및 권고 중심의 프레임워크로서, 국가 정책과 기업 전략 수립의 방향성을 제시한다. 둘째, ISO/IEC, IEEE, NIST가 대표하는 표준 및 프레임워크 중심의 접근으로서, 조직 수준의 실행 가이드를 제공한다. 실무적으로는 양자를 함께 참조하여 거버넌스 체계를 설계하는 것이 일반적이다.

%----------------------------------------------------------
\subsection{OECD AI 원칙(OECD AI Principles)}
%----------------------------------------------------------
OECD(경제협력개발기구)는 2019년 5월 ``AI에 관한 OECD 권고''를 채택하였으며\cite{oecd2019principles}, 이는 2024년 5월 생성형 AI 및 범용 AI의 등장에 대응하여 개정되었다\cite{oecd2024framework}. 이는 글로벌 AI 정책의 기준점으로 기능한다.

\subsubsection{5대 가치 기반 원칙 심층 분석}

OECD AI 원칙은 다섯 가지 가치 기반 원칙과 다섯 가지 정책 권고로 구성된다. 각 원칙의 실무적 의미와 적용 사례를 상세히 분석한다.

\paragraph{원칙 1: 포용적 성장, 지속가능한 발전, 웰빙}
AI가 인간과 지구의 웰빙에 기여해야 한다는 원칙이다. 이는 단순히 경제적 이익을 넘어, AI 기술이 사회 전반의 복지를 향상시키고 환경적 지속가능성에 기여해야 함을 의미한다.

\begin{itemize}
 \item \textbf{실무 적용}: AI 프로젝트 기획 시 사회적 영향 평가(Social Impact Assessment)를 수행하여, 해당 AI 시스템이 특정 집단의 소외를 야기하지 않는지 검증한다.
 \item \textbf{적용 사례}: 한국의 한 금융기관은 AI 신용평가 모델 도입 시, 금융 소외 계층(저소득층, 고령자, 신용 이력 부족자)에 대한 접근성 향상 지표를 핵심 성과 지표(KPI)로 설정하였다.
 \item \textbf{환경적 고려}: AI 모델 학습에 소요되는 에너지 소비량을 측정하고, 탄소 배출 저감 목표를 설정하는 것도 이 원칙의 실무적 구현에 해당한다.
\end{itemize}

\paragraph{원칙 2: 법치, 인권, 민주적 가치 존중}
공정성(Fairness), 프라이버시(Privacy), 비차별(Non-discrimination)을 보장해야 한다는 원칙이다.

\begin{itemize}
 \item \textbf{실무 적용}: AI 시스템의 입력 데이터와 출력 결과에 대한 공정성 감사(Fairness Audit)를 정기적으로 수행한다. 보호 속성(성별, 연령, 인종, 장애 여부 등)에 따른 차별적 결과가 발생하지 않는지 통계적으로 검증한다.
 \item \textbf{적용 사례}: 채용 AI를 운영하는 기업이 인구통계학적 동등성(Demographic Parity), 균등 기회(Equalized Odds), 예측 평등(Predictive Equality) 등 다중 공정성 지표를 모니터링하는 체계를 구축한 사례가 있다.
 \item \textbf{프라이버시 보호}: 개인정보 처리 최소화 원칙 적용, 차분 프라이버시(Differential Privacy) 기법 도입, 연합학습(Federated Learning) 활용 등이 이 원칙의 기술적 구현 방안이다.
\end{itemize}

\paragraph{원칙 3: 투명성 및 설명가능성}
AI 시스템의 작동 방식에 대한 이해 및 이의 제기 가능성을 보장해야 한다.\cite{arrieta2020explainable}

\begin{itemize}
 \item \textbf{실무 적용}: 모델 카드(Model Card) 작성을 의무화하고, 주요 의사결정에 대한 설명 보고서(Explanation Report)를 생성하는 시스템을 구축한다.
 \item \textbf{적용 사례}: 보험사의 AI 심사 시스템에서 SHAP(SHapley Additive exPlanations) 값을 기반으로 개별 심사 결과에 대한 요인 분석 보고서를 자동 생성하여 고객에게 제공하는 체계를 도입한 사례가 있다.
 \item \textbf{기술적 구현}: LIME(Local Interpretable Model-agnostic Explanations), SHAP, Attention 시각화 등 설명가능성 기법을 AI 시스템에 내장하고, 이해관계자 유형별(기술자, 경영진, 규제기관, 최종 사용자) 맞춤형 설명을 제공한다.
\end{itemize}

\paragraph{원칙 4: 견고성, 보안, 안전}
정상 및 비정상 조건에서의 안전한 작동을 보장해야 한다.

\begin{itemize}
 \item \textbf{실무 적용}: 적대적 공격(Adversarial Attack) 테스트, 스트레스 테스트, 장애 복구(Failover) 시나리오 테스트를 정기적으로 수행한다.
 \item \textbf{적용 사례}: 자율주행 AI를 개발하는 기업이 100만 건 이상의 에지 케이스(Edge Case) 시나리오를 시뮬레이션하고, 안전 임계값 위반 시 자동으로 인간 제어 모드로 전환되는 안전장치(Fail-safe)를 구현한 사례가 있다.
 \item \textbf{보안 관점}: 모델 도난(Model Stealing), 데이터 포이즈닝(Data Poisoning), 프롬프트 인젝션(Prompt Injection) 등 AI 특유의 보안 위협에 대한 방어 체계를 수립한다.
\end{itemize}

\paragraph{원칙 5: 책임성}
AI 행위자의 적절한 기능과 원칙 준수에 대한 책무를 명확히 해야 한다.

\begin{itemize}
 \item \textbf{실무 적용}: AI 시스템의 전체 생명주기에 걸쳐 역할과 책임을 명확히 정의하는 RACI(Responsible, Accountable, Consulted, Informed) 매트릭스를 수립한다.
 \item \textbf{적용 사례}: AI 시스템의 의사결정으로 인한 피해 발생 시, 피해자 구제 절차와 이의신청 채널을 사전에 설계하고 운영하는 기업 사례가 있다.
 \item \textbf{거버넌스 체계}: AI 윤리위원회, AI 리스크 관리위원회 등 의사결정 기구를 설치하고, 정기적 보고 체계를 운영한다.
\end{itemize}

\subsubsection{2024년 개정 상세 분석}
2024년 개정은 급격히 발전한 생성형 AI(Generative AI)와 범용 AI(General-Purpose AI)가 야기하는 새로운 리스크에 대응하기 위한 것이다. 주요 개정 사항을 상세히 분석한다.

\begin{itemize}
 \item \textbf{정보 무결성(Information Integrity) 강화}: 생성형 AI로 인한 허위정보(Disinformation) 확산에 대응하기 위해, AI 생성 콘텐츠의 출처 표시와 진위 확인 메커니즘을 강조한다. 이는 C2PA(Coalition for Content Provenance and Authenticity) 표준과의 연계를 권고한다.
 \item \textbf{안전 관련 조항 구체화}: 의도하지 않은 해악 방지를 위한 견고한 메커니즘과 인간 제어권(Human Oversight)을 강조한다. 특히 고위험 AI 시스템에 대한 ``킬 스위치(Kill Switch)'' 메커니즘과 인간 개입(Human-in-the-Loop) 요건을 구체화하였다.
 \item \textbf{책임성 원칙 강화}: 추적가능성(Traceability), 체계적 리스크관리(Systematic Risk Management), 지적재산권(Intellectual Property) 고려 등을 명시하였다. AI 시스템의 학습 데이터 출처와 모델 변경 이력을 기록하는 감사 추적(Audit Trail) 요건이 추가되었다.
 \item \textbf{환경 영향 고려}: AI 시스템의 에너지 소비와 환경 영향에 대한 고려가 명시적으로 추가되었다. 대규모 언어 모델(LLM) 학습의 탄소 발자국(Carbon Footprint) 측정 및 보고를 권장한다.
 \item \textbf{다중이해관계자 참여}: AI 시스템의 설계와 배포 과정에서 영향을 받는 집단의 의견을 수렴하는 절차를 강화하였다.
\end{itemize}

\subsubsection{한국 기업 적용 체크리스트}

OECD AI 원칙은 연성 규범(Soft Law)이지만, EU AI Act, 미국의 AI 정책 등 주요국 규제의 기반이 된다. 다국적 사업을 영위하는 한국 기업이 이를 자사 거버넌스 체계의 기반으로 삼는 것은 글로벌 규제 환경에 선제적으로 대응하는 전략이 된다.

\begin{practicaltool}{OECD AI 원칙 기반 한국 기업 자가진단 체크리스트}
다음 항목을 분기별로 점검하여 OECD AI 원칙 준수 수준을 진단한다.

\begin{center}
\renewcommand{\arraystretch}{1.3}
\small
\begin{tabularx}{\textwidth}{p{0.25\textwidth}Xp{0.1\textwidth}}
\toprule
\textbf{원칙 영역} & \textbf{점검 항목} & \textbf{준수 여부} \\
\midrule
포용적 성장 & AI 프로젝트의 사회적 영향 평가를 수행하고 있는가? & $\square$ \\
 & 취약 계층에 대한 접근성을 고려하고 있는가? & $\square$ \\
 & AI 학습의 환경 영향(에너지 소비)을 측정하고 있는가? & $\square$ \\
\midrule
인권·민주적 가치 & 공정성 감사를 정기적으로 수행하고 있는가? & $\square$ \\
 & 개인정보보호 영향평가(PIA)를 수행하고 있는가? & $\square$ \\
 & 비차별 정책이 문서화되어 있는가? & $\square$ \\
\midrule
투명성·설명가능성 & 모델 카드(Model Card)를 작성하고 있는가? & $\square$ \\
 & 이해관계자별 맞춤 설명을 제공하고 있는가? & $\square$ \\
 & 이의신청 채널이 운영되고 있는가? & $\square$ \\
\midrule
견고성·보안·안전 & 적대적 공격 테스트를 수행하고 있는가? & $\square$ \\
 & 장애 복구 시나리오가 설계되어 있는가? & $\square$ \\
 & AI 보안 위협(프롬프트 인젝션 등) 대응 체계가 있는가? & $\square$ \\
\midrule
책임성 & RACI 매트릭스가 정의되어 있는가? & $\square$ \\
 & 감사 추적(Audit Trail)이 기록되고 있는가? & $\square$ \\
 & 피해 구제 절차가 수립되어 있는가? & $\square$ \\
\bottomrule
\end{tabularx}
\end{center}
\end{practicaltool}

%----------------------------------------------------------
\subsection{UNESCO AI 윤리 권고안}
%----------------------------------------------------------
2021년 11월 채택된 UNESCO 권고안은 193개 회원국의 만장일치로 통과된 최초의 글로벌 규범 문서로, 보편성을 갖는다.\cite{unesco2021ethics}

\subsubsection*{10대 원칙}
\begin{enumerate}
 \item 비례성 및 해악 금지
 \item 안전 및 보안
 \item 공정성 및 비차별
 \item 지속가능성
 \item 프라이버시권 및 데이터 보호
 \item 인간 감독 및 결정권
 \item 투명성 및 설명가능성
 \item 책임 및 책무성
 \item 인식 및 리터러시
 \item 다중이해관계자 및 적응적 거버넌스
\end{enumerate}

\subsubsection*{이행 도구: RAM과 EIA}
UNESCO는 실질적인 이행 도구를 제공한다.
\begin{itemize}
 \item \textbf{준비도 평가 방법론(Readiness Assessment Methodology, RAM)}: 국가 및 조직의 AI 윤리 준비도를 평가하기 위한 정량적·정성적 지표 체계이다. 법·제도, 기술 인프라, 인적 역량, 사회적 인식의 네 가지 차원에서 평가한다.
 \item \textbf{윤리적 영향 평가(Ethical Impact Assessment, EIA)}: 개별 프로젝트 수준의 잠재적 영향 평가 도구이다. AI 시스템이 인권, 사회적 웰빙, 환경에 미치는 영향을 사전에 식별하고 완화 방안을 수립한다.
\end{itemize}
SAP 등 글로벌 기업들이 자사 AI 윤리 정책을 UNESCO 권고안에 정렬(Alignment)하는 사례가 증가하고 있다.

%----------------------------------------------------------
\subsection{ISO/IEC 42001 AI 관리체계 표준}
%----------------------------------------------------------
ISO/IEC 42001:2023은 세계 최초의 AI 관리체계(AI Management System, AIMS) 국제표준이다.\cite{iso42001} ISO의 고수준 구조(Harmonized Structure, HS)를 따르며 PDCA(Plan-Do-Check-Act) 방법론에 기반한다.

\subsubsection{10대 조항 구조 상세 분석}

ISO/IEC 42001은 ISO 경영시스템 표준의 공통 구조인 10개 조항으로 구성된다. 각 조항의 요구사항과 AI 거버넌스에서의 실무적 의미를 상세히 분석한다.

\begin{table}[htbp]
\centering
\caption{ISO/IEC 42001 핵심 요구사항}
\label{tab:iso42001_requirements}
\begin{tabularx}{\textwidth}{p{0.5cm}p{2.5cm}X}
\toprule
\textbf{No.} & \textbf{조항명} & \textbf{주요 요구사항} \\
\midrule
1 & 적용 범위 & 표준의 적용 범위와 목적 정의 \\
2 & 인용 표준 & 참조하는 관련 표준 목록 \\
3 & 용어 및 정의 & AI 관리체계 관련 핵심 용어 정의 \\
4 & 조직의 맥락 & 내외부 이슈 파악, 이해관계자 요구 분석, AIMS 범위 결정 \\
5 & 리더십 & 최고경영진의 의지 표명, AI 정책 수립, 역할·책임·권한 배정 \\
6 & 기획 & 리스크 및 기회 식별, AI 목표 수립, 변경 계획 \\
7 & 지원 & 자원 확보, 역량 개발, 인식 제고, 의사소통, 문서화된 정보 관리 \\
8 & 운영 & 운영 기획 및 제어, AI 영향 평가, AI 시스템 생명주기 관리, 데이터 관리, 제3자 관리 \\
9 & 성능 평가 & 모니터링·측정·분석·평가, 내부 감사, 경영 검토 \\
10 & 개선 & 부적합 및 시정조치, 지속적 개선 \\
\bottomrule
\end{tabularx}
\end{table}

\paragraph{조항 4: 조직의 맥락}
조직은 AI 관리체계에 관련된 내부 및 외부 이슈를 결정해야 한다. 내부 이슈에는 조직의 AI 역량, 문화, 기술 인프라가 포함되며, 외부 이슈에는 규제 환경, 시장 동향, 기술 발전이 포함된다. 특히 이해관계자의 요구사항과 기대를 파악하고, 이를 AIMS의 범위에 반영해야 한다.

\paragraph{조항 5: 리더십}
최고경영진은 AI 관리체계에 대한 리더십과 의지를 표명해야 한다. 이는 AI 정책의 수립과 공표, 역할과 책임의 명확한 배정, 필요한 자원의 할당을 포함한다. AI 정책은 조직의 목적에 적합하고, 목표 설정의 프레임워크를 제공하며, 적용 가능한 요구사항의 충족에 대한 의지를 포함해야 한다.

\paragraph{조항 6: 기획}
조직은 AI 시스템과 관련된 리스크와 기회를 식별하고 이에 대응하는 조치를 계획해야 한다. 여기에는 AI 영향 평가(AI Impact Assessment)의 수행, AI 시스템의 리스크 분류, 잔류 리스크(Residual Risk)의 수용 기준 설정이 포함된다.

\paragraph{조항 8: 운영}
운영 조항은 ISO/IEC 42001의 핵심이다. AI 시스템 생명주기 전체에 걸친 관리 절차를 요구하며, 다음의 세부 영역을 다룬다:
\begin{itemize}
 \item \textbf{AI 영향 평가}: AI 시스템이 개인, 집단, 사회에 미치는 영향을 체계적으로 평가
 \item \textbf{AI 시스템 생명주기}: 설계, 개발, 검증, 배포, 운영, 폐기에 이르는 전 과정 관리
 \item \textbf{데이터 관리}: 학습 데이터, 검증 데이터, 운영 데이터의 품질, 편향, 프라이버시 관리
 \item \textbf{제3자 및 공급업체 관리}: 외부에서 조달하는 AI 모델, 데이터, 서비스에 대한 관리
\end{itemize}

\subsubsection{Annex A 제어 목록 주요 항목}

Annex A는 조직이 선택적으로 적용할 수 있는 AI 특화 제어 항목을 제공한다. 주요 영역은 다음과 같다:

\begin{itemize}
 \item \textbf{A.2 AI 정책}: AI 사용에 관한 내부 정책 수립 및 이해관계자 공개
 \item \textbf{A.3 내부 조직}: AI 거버넌스 역할 및 책임 정의, AI 시스템 인벤토리 관리
 \item \textbf{A.4 AI 시스템 자원}: 데이터 관리, 도구 및 프레임워크 관리, 시스템 및 컴퓨팅 자원 관리
 \item \textbf{A.5 AI 시스템 영향 평가}: 영향 평가 수행 절차, 평가 결과의 문서화 및 검토
 \item \textbf{A.6 AI 시스템 생명주기}: 요구사항 정의, 설계 및 개발, 검증 및 확인, 배포, 운영 및 모니터링
 \item \textbf{A.7 데이터}: 데이터 출처 관리, 데이터 품질, 데이터 준비 및 라벨링
 \item \textbf{A.8 정보 및 문서화}: AI 시스템 문서화, 모델 설명, 이해관계자 커뮤니케이션
 \item \textbf{A.9 제3자 및 고객 관계}: 공급업체 관리, 고객 커뮤니케이션, 제3자 AI 구성요소 관리
\end{itemize}

\subsubsection{인증 프로세스}

ISO/IEC 42001 인증은 다음 네 단계로 진행된다:

\paragraph{1단계: 인증 준비}
\begin{itemize}
 \item AIMS 범위 정의 및 적용 가능성 분석(Gap Analysis) 수행
 \item AI 정책, 리스크 평가 방법론, 영향 평가 절차 등 필수 문서 수립
 \item 내부 감사원 교육 및 인증 기관 선정
 \item 준비 기간: 일반적으로 6개월$\sim$12개월
\end{itemize}

\paragraph{2단계: 인증 심사}
\begin{itemize}
 \item \textbf{1단계 심사(Stage 1 Audit)}: 문서 검토 중심. AIMS의 설계 적절성을 평가한다. 필수 문서(AI 정책, 리스크 평가 결과, 영향 평가 결과, 적용성 보고서 등)의 존재와 적합성을 확인한다.
 \item \textbf{2단계 심사(Stage 2 Audit)}: 현장 심사 중심. AIMS의 구현 효과성을 평가한다. 실제 AI 시스템의 운영 현황, 기록, 인터뷰를 통해 문서화된 절차가 실행되고 있는지 확인한다.
\end{itemize}

\paragraph{3단계: 인증 발급}
\begin{itemize}
 \item 심사 결과에 따라 적합성 판정 및 인증서 발급
 \item 부적합(Non-Conformity) 사항이 발견된 경우, 시정조치 완료 후 인증 발급
 \item 인증 유효기간: 3년
\end{itemize}

\paragraph{4단계: 사후관리}
\begin{itemize}
 \item \textbf{사후 심사(Surveillance Audit)}: 연 1회 이상 수행. AIMS의 지속적 적합성 확인
 \item \textbf{갱신 심사(Recertification Audit)}: 3년 주기. 전체 AIMS에 대한 재평가
 \item 중대한 변경(AI 시스템 범위 확대, 조직 구조 변경 등) 발생 시 추가 심사 가능
\end{itemize}

\subsubsection{ISO 42001 인증 현황 (2026년 기준)}

\begin{tcolorbox}[colback=blue!5, colframe=blue!60, title={ISO 42001 인증 시장 현황 (2026년 초 기준)}]
2026년 초 기준, 전 세계 ISO/IEC 42001 인증 취득 기업은 약 \textbf{16개사}에 불과하다. Microsoft, KPMG, IBM, Synthesia 등 주요 기업이 포함되어 있으나, 감사 기관의 역량이 수요를 따라가지 못하는 상황이다. 그러나 시장은 ``첫 번째 실질적 성장 국면''에 진입하였다:
\begin{itemize}
    \item 기업의 \textbf{76\%}가 향후 24개월 내 AI 감사 또는 인증을 추진할 계획
    \item ISO 42001 준수가 EU AI Act 고위험 시스템 문서화 요건의 \textbf{약 70\%}를 커버
    \item 인증 취득 기업은 EU 시장에서 규제 준수 입증 비용을 크게 절감
\end{itemize}
\end{tcolorbox}

\subsubsection{타 ISO 표준과의 통합 방안}

ISO/IEC 42001은 ISO 고수준 구조(HS)를 따르므로, 기존에 ISO 27001(정보보안), ISO 9001(품질경영) 등을 운영하는 조직은 통합 관리체계(Integrated Management System, IMS)를 구축할 수 있다.

\begin{itemize}
 \item \textbf{ISO 27001과의 통합}: AI 시스템의 정보보안 측면을 ISO 27001 체계로 관리하고, AI 특화 리스크(모델 보안, 데이터 포이즈닝 등)은 ISO 42001로 관리한다. 공통 영역(리스크 평가, 내부 감사, 경영 검토)은 통합 운영한다.
 \item \textbf{ISO 9001과의 통합}: AI 시스템의 품질 관리(모델 성능, 데이터 품질)를 ISO 9001의 품질경영 체계와 연계한다. 프로세스 접근법(Process Approach)과 지속적 개선(Continual Improvement)의 원칙을 공유한다.
 \item \textbf{통합의 이점}: 중복 문서화 감소, 심사 효율성 향상, 조직 전체의 관리체계 일관성 확보, 비용 절감 효과가 있다.
\end{itemize}

%----------------------------------------------------------
\subsection{IEEE 7000 시리즈}
%----------------------------------------------------------
IEEE 7000-2021은 시스템 설계 과정에서 윤리적 문제를 식별하고 분석하기 위한 \textbf{가치 기반 공학(Value-Based Engineering)} 접근법을 제시한다.\cite{ieee2019ead}

\subsubsection*{주요 프로세스}
\begin{itemize}
 \item \textbf{맥락 탐색 및 가치 도출}: 이해관계자로부터 윤리적 가치 도출.
 \item \textbf{가치 요구사항 공식화}: 추상적 가치를 구체적 시스템 요구사항으로 변환.
 \item \textbf{가치 등록부(Value Register)}: 가치와 요구사항의 추적성 기록.
\end{itemize}
IEEE 7001(투명성), 7002(프라이버시), 7003(알고리즘 편향) 등 확장된 표준군과 \textbf{CertifAIEd} 인증 프로그램을 운영 중이다.

%----------------------------------------------------------
\subsection{NIST AI 리스크관리 프레임워크(AI RMF)}
%----------------------------------------------------------
미국 NIST가 발표한 AI RMF 1.0은 AI 시스템의 신뢰성 고려를 통합하기 위한 자발적 가이드를 제공한다.\cite{nist2023airmf}

\subsubsection{4대 핵심 기능 상세 분석}

NIST AI RMF는 거버넌스(Govern), 매핑(Map), 측정(Measure), 관리(Manage)의 4대 핵심 기능으로 구성되며, 각 기능은 카테고리와 하위 카테고리로 세분화된다.

\paragraph{Govern (거버넌스)}
조직 문화, 정책, 프로세스, 절차를 수립하여 AI 리스크관리를 위한 기반을 마련하는 기능이다. 다른 세 가지 기능 전체에 걸쳐 적용되는 교차(Cross-cutting) 기능이다.

\begin{itemize}
 \item \textbf{GV-1}: AI 시스템에 대한 법적·규제적 요구사항을 식별하고 준수한다.
 \item \textbf{GV-2}: AI 리스크관리에 관한 책임성 구조와 프로세스를 수립한다.
 \item \textbf{GV-3}: 인력의 다양성(Diversity), 형평성(Equity), 포용성(Inclusion), 접근성(Accessibility)을 조직 정책에 반영한다.
 \item \textbf{GV-4}: 조직의 팀이 AI 리스크를 인식하고, 이를 평가·대응할 역량을 갖추도록 한다.
 \item \textbf{GV-5}: AI 시스템에 대한 강건한 참여 메커니즘(Engagement Mechanism)을 수립한다.
 \item \textbf{GV-6}: AI 리스크관리 정책 및 프로세스를 문서화하고 이를 조직 내 공유한다.
\end{itemize}

\paragraph{Map (매핑)}
AI 시스템이 사용되는 맥락을 이해하고, 관련 리스크를 식별하는 기능이다.

\begin{itemize}
 \item \textbf{MP-1}: AI 시스템의 의도된 목적, 잠재적 유익한 사용, 유해한 오용을 식별한다.
 \item \textbf{MP-2}: AI 시스템의 과학적·기술적 기반에 대한 이해를 확보한다.
 \item \textbf{MP-3}: AI 시스템의 맥락에 관련된 리스크와 이점을 체계적으로 식별한다.
 \item \textbf{MP-4}: AI 시스템과 관련된 리스크를 우선순위화하고, 리스크 허용 범위(Risk Tolerance)를 설정한다.
 \item \textbf{MP-5}: AI 시스템에 영향을 받는 이해관계자를 식별하고 참여시킨다.
\end{itemize}

\paragraph{Measure (측정)}
식별된 리스크를 정량적·정성적으로 분석하고 평가하는 기능이다.

\begin{itemize}
 \item \textbf{ME-1}: AI 시스템과 관련된 리스크를 측정하기 위한 적절한 방법론과 지표를 선정한다.
 \item \textbf{ME-2}: AI 시스템을 테스트하고, 그 결과를 분석·보고한다.
 \item \textbf{ME-3}: AI 시스템의 신뢰성 특성(정확성, 공정성, 견고성, 프라이버시 등)을 지속적으로 모니터링한다.
 \item \textbf{ME-4}: 측정 결과를 검증하고, 측정 방법론의 적절성을 평가한다.
\end{itemize}

\paragraph{Manage (관리)}
식별·측정된 리스크에 대응하고, 이를 지속적으로 모니터링하는 기능이다.

\begin{itemize}
 \item \textbf{MG-1}: 리스크 대응 계획을 수립하고 우선순위화한다.
 \item \textbf{MG-2}: 리스크 대응 조치를 이행하고 그 효과를 추적한다.
 \item \textbf{MG-3}: 사전에 정의된 리스크 관리 계획과 관련 결정에 대한 이의를 제기할 수 있는 절차를 마련한다.
 \item \textbf{MG-4}: AI 리스크와 관련 의사결정에 대한 리스크 의사소통(Risk Communication)을 수행한다.
\end{itemize}

\subsubsection{생성형 AI 프로파일: 12대 리스크 상세}

2024년 7월 발표된 NIST AI 600-1 ``생성형 AI 프로파일(Generative AI Profile)''은 생성형 AI 고유의 리스크에 대한 대응을 구체화한 문서이다. 12대 리스크 유형을 상세히 분석한다.

\begin{enumerate}
 \item \textbf{CBRN 정보 접근(CBRN Information)}: 화학·생물·방사능·핵 무기 관련 위험한 정보의 생성 가능성
 \item \textbf{허위 콘텐츠(Confabulation)}: 사실과 다른 정보를 자신 있게 생성하는 ``환각(Hallucination)'' 현상
 \item \textbf{데이터 프라이버시(Data Privacy)}: 학습 데이터에 포함된 개인정보의 노출 리스크
 \item \textbf{환경 영향(Environmental Impact)}: 대규모 모델 학습 및 추론의 에너지 소비와 탄소 배출
 \item \textbf{유해 편향 및 동질화(Harmful Bias and Homogenization)}: 편향된 출력 생성 및 다양성 감소
 \item \textbf{인간-AI 상호작용(Human-AI Configuration)}: 과도한 의존(Over-reliance), 자동화 편향(Automation Bias)
 \item \textbf{정보 무결성(Information Integrity)}: 허위정보 생성 및 확산, 여론 조작 가능성
 \item \textbf{정보 보안(Information Security)}: 프롬프트 인젝션, 간접 프롬프트 인젝션, 모델 탈옥(Jailbreaking)
 \item \textbf{지적재산권(Intellectual Property)}: 학습 데이터의 저작권 침해, 생성물의 저작권 문제
 \item \textbf{불투명성·설명가능성(Obscene, Degrading Content)}: 유해하거나 부적절한 콘텐츠 생성
 \item \textbf{가치 연계(Value Chain and Component Integration)}: 제3자 모델·데이터·서비스 통합 시 리스크 전파
 \item \textbf{악의적 사용(Malicious Use)}: 사이버 공격, 사기, 사회공학적 공격에의 악용
\end{enumerate}

\subsubsection{NIST AI RMF 실무 적용 프로세스}

NIST AI RMF를 조직에 적용하기 위한 실무 프로세스는 다음과 같다:

\begin{enumerate}
 \item \textbf{조직 맥락 파악}: 조직의 AI 활용 현황, 규제 환경, 이해관계자 기대를 파악한다.
 \item \textbf{AI 시스템 인벤토리}: 조직 내 모든 AI 시스템을 식별하고 목록화한다.
 \item \textbf{리스크 프로파일 작성}: 각 AI 시스템에 대해 4대 기능별 리스크 프로파일을 작성한다.
 \item \textbf{격차 분석}: 현재 리스크관리 수준과 목표 수준 간의 격차를 파악한다.
 \item \textbf{대응 계획 수립}: 격차를 해소하기 위한 구체적 조치 계획을 수립한다.
 \item \textbf{이행 및 모니터링}: 계획을 이행하고, 효과를 지속적으로 모니터링한다.
\end{enumerate}

\begin{practicaltool}{NIST AI RMF 준수 자동 평가 도구}
다음 Python 코드는 NIST AI RMF의 4대 기능별 준수 상태를 자동으로 평가하고 보고서를 생성하는 도구이다.

\begin{lstlisting}[language=Python]
"""NIST AI RMF Compliance Auto-Assessment Tool"""
from dataclasses import dataclass, field
from typing import Dict, List, Optional
from enum import Enum
import json
from datetime import datetime

class ComplianceLevel(Enum):
    NOT_STARTED = 0
    PARTIAL = 1
    SUBSTANTIAL = 2
    FULL = 3

class RMFFunction(Enum):
    GOVERN = "Govern"
    MAP = "Map"
    MEASURE = "Measure"
    MANAGE = "Manage"

@dataclass
class SubCategory:
    id: str
    description: str
    compliance: ComplianceLevel = ComplianceLevel.NOT_STARTED
    evidence: str = ""
    action_plan: str = ""

@dataclass
class Category:
    id: str
    name: str
    function: RMFFunction
    subcategories: List[SubCategory] = field(
        default_factory=list
    )

    @property
    def compliance_score(self) -> float:
        if not self.subcategories:
            return 0.0
        total = sum(
            sc.compliance.value
            for sc in self.subcategories
        )
        return total / (
            len(self.subcategories)
            * ComplianceLevel.FULL.value
        ) * 100

class NISTAIRMFAssessment:
    """NIST AI RMF compliance assessment framework."""

    def __init__(self, org_name: str, system_name: str):
        self.org_name = org_name
        self.system_name = system_name
        self.assessment_date = datetime.now()
        self.categories: Dict[str, Category] = {}
        self._initialize_framework()

    def _initialize_framework(self):
        """Initialize all RMF categories and subcategories."""
        govern_cats = {
            "GV-1": ("Legal & Regulatory",
                      ["GV-1.1", "GV-1.2"]),
            "GV-2": ("Accountability Structure",
                      ["GV-2.1", "GV-2.2"]),
            "GV-3": ("Workforce Diversity",
                      ["GV-3.1", "GV-3.2"]),
            "GV-4": ("Org Competency",
                      ["GV-4.1", "GV-4.2", "GV-4.3"]),
            "GV-5": ("Engagement Mechanisms",
                      ["GV-5.1", "GV-5.2"]),
            "GV-6": ("Policies Documented",
                      ["GV-6.1", "GV-6.2"]),
        }
        for cat_id, (name, subs) in govern_cats.items():
            self.categories[cat_id] = Category(
                id=cat_id,
                name=name,
                function=RMFFunction.GOVERN,
                subcategories=[
                    SubCategory(id=s, description=f"{name} - {s}")
                    for s in subs
                ],
            )

        map_cats = {
            "MP-1": ("Intended Purpose", ["MP-1.1", "MP-1.2"]),
            "MP-2": ("Scientific Understanding",
                      ["MP-2.1", "MP-2.2"]),
            "MP-3": ("Risk Identification",
                      ["MP-3.1", "MP-3.2"]),
            "MP-4": ("Risk Prioritization", ["MP-4.1"]),
            "MP-5": ("Stakeholder Engagement",
                      ["MP-5.1", "MP-5.2"]),
        }
        for cat_id, (name, subs) in map_cats.items():
            self.categories[cat_id] = Category(
                id=cat_id,
                name=name,
                function=RMFFunction.MAP,
                subcategories=[
                    SubCategory(id=s, description=f"{name} - {s}")
                    for s in subs
                ],
            )

        measure_cats = {
            "ME-1": ("Metrics Selection", ["ME-1.1", "ME-1.2"]),
            "ME-2": ("Testing & Analysis", ["ME-2.1", "ME-2.2"]),
            "ME-3": ("Continuous Monitoring",
                      ["ME-3.1", "ME-3.2"]),
            "ME-4": ("Measurement Validation", ["ME-4.1"]),
        }
        for cat_id, (name, subs) in measure_cats.items():
            self.categories[cat_id] = Category(
                id=cat_id,
                name=name,
                function=RMFFunction.MEASURE,
                subcategories=[
                    SubCategory(id=s, description=f"{name} - {s}")
                    for s in subs
                ],
            )

        manage_cats = {
            "MG-1": ("Response Planning", ["MG-1.1", "MG-1.2"]),
            "MG-2": ("Response Implementation",
                      ["MG-2.1", "MG-2.2"]),
            "MG-3": ("Appeal Mechanisms", ["MG-3.1"]),
            "MG-4": ("Risk Communication", ["MG-4.1", "MG-4.2"]),
        }
        for cat_id, (name, subs) in manage_cats.items():
            self.categories[cat_id] = Category(
                id=cat_id,
                name=name,
                function=RMFFunction.MANAGE,
                subcategories=[
                    SubCategory(id=s, description=f"{name} - {s}")
                    for s in subs
                ],
            )

    def update_assessment(
        self,
        subcategory_id: str,
        compliance: ComplianceLevel,
        evidence: str = "",
        action_plan: str = "",
    ):
        """Update compliance status for a subcategory."""
        for cat in self.categories.values():
            for sc in cat.subcategories:
                if sc.id == subcategory_id:
                    sc.compliance = compliance
                    sc.evidence = evidence
                    sc.action_plan = action_plan
                    return
        raise ValueError(
            f"Subcategory {subcategory_id} not found"
        )

    def get_function_score(self, func: RMFFunction) -> float:
        """Calculate overall score for an RMF function."""
        func_cats = [
            c for c in self.categories.values()
            if c.function == func
        ]
        if not func_cats:
            return 0.0
        return sum(
            c.compliance_score for c in func_cats
        ) / len(func_cats)

    def get_overall_score(self) -> float:
        """Calculate overall compliance score."""
        return sum(
            self.get_function_score(f)
            for f in RMFFunction
        ) / len(RMFFunction)

    def generate_report(self) -> dict:
        """Generate compliance assessment report."""
        report = {
            "organization": self.org_name,
            "system": self.system_name,
            "assessment_date": self.assessment_date.isoformat(),
            "overall_score": round(
                self.get_overall_score(), 1
            ),
            "function_scores": {},
            "gap_analysis": [],
        }
        for func in RMFFunction:
            score = self.get_function_score(func)
            report["function_scores"][func.value] = round(
                score, 1
            )

        # Identify gaps (below substantial compliance)
        for cat in self.categories.values():
            for sc in cat.subcategories:
                if sc.compliance.value < ComplianceLevel.SUBSTANTIAL.value:
                    report["gap_analysis"].append({
                        "subcategory": sc.id,
                        "current_level": sc.compliance.name,
                        "description": sc.description,
                        "action_plan": sc.action_plan,
                    })
        return report


# Usage example
if __name__ == "__main__":
    assessment = NISTAIRMFAssessment(
        org_name="Example Corp Korea",
        system_name="AI Credit Scoring v2.1",
    )

    # Assess several subcategories
    assessment.update_assessment(
        "GV-1.1",
        ComplianceLevel.FULL,
        evidence="Regulatory tracking system deployed",
    )
    assessment.update_assessment(
        "GV-2.1",
        ComplianceLevel.SUBSTANTIAL,
        evidence="RACI matrix documented",
    )
    assessment.update_assessment(
        "MP-3.1",
        ComplianceLevel.PARTIAL,
        evidence="Initial risk register created",
        action_plan="Complete risk assessment by Q2",
    )
    assessment.update_assessment(
        "ME-1.1",
        ComplianceLevel.NOT_STARTED,
        action_plan="Define fairness metrics by Q1",
    )

    report = assessment.generate_report()
    print(json.dumps(report, indent=2, ensure_ascii=False))
    print(f"\nOverall Score: {report['overall_score']}%")
    print(f"Gaps found: {len(report['gap_analysis'])}")
\end{lstlisting}
\end{practicaltool}

%----------------------------------------------------------
\subsection{프레임워크 간 상호운용성 정렬}
%----------------------------------------------------------
대부분의 전문가들은 다음과 같은 통합적 접근을 권장한다:
\begin{itemize}
 \item \textbf{원칙 수준}: OECD/UNESCO (윤리 원칙 채택)
 \item \textbf{리스크관리 수준}: NIST AI RMF (리스크 식별/평가)
 \item \textbf{관리체계 수준}: ISO/IEC 42001 (공식 체계 구축 및 인증)
 \item \textbf{설계 수준}: IEEE 7000 (설계 요구사항 변환)
\end{itemize}

이러한 계층적 접근은 각 프레임워크의 강점을 활용하면서 중복을 최소화한다. 실무적으로는 ISO/IEC 42001을 중심축으로 삼고, NIST AI RMF의 리스크관리 방법론을 리스크 평가 절차에 통합하며, OECD/UNESCO 원칙을 AI 정책의 기반으로 활용하는 방식이 효과적이다.

%==========================================================
\section{주요국 규제 동향}
%==========================================================
AI 거버넌스의 국제적 지형은 2026년 1월 현재, EU의 리스크 기반 규제와 미국의 트럼프 행정부 주도 자율 규제 전환이라는 두 축을 중심으로 전개되고 있다.\cite{smuha2021race, bradford2020brussels}

%----------------------------------------------------------
\subsection{유럽연합: EU AI Act 심층 분석}
%----------------------------------------------------------
2024년 8월 발효된 EU AI Act는 세계 최초의 포괄적 AI 규제 법률로서 글로벌 표준으로 자리 잡고 있다.\cite{eu2024aiact} 이는 전 세계 100개 이상의 AI 윤리 가이드라인 분석 결과와도 그 궤를 같이한다.\cite{jobin2019global}

\subsubsection{4단계 리스크 분류 체계}

EU AI Act의 핵심은 AI 시스템을 리스크 수준에 따라 4단계로 분류하는 리스크 기반 접근법이다.

\paragraph{1단계: 금지 AI 시스템(Unacceptable Risk)}
2025년 2월 2일부터 시행된 금지 규정은 EU의 기본권과 가치에 반하는 AI 활용을 원천적으로 금지한다. 금지 대상은 다음과 같다:

\begin{enumerate}
 \item \textbf{잠재의식 조작(Subliminal Manipulation)}: 인간이 인지하지 못하는 방식으로 행동을 조작하여 해를 끼치는 AI 시스템
 \item \textbf{취약 계층 착취(Exploitation of Vulnerabilities)}: 연령, 장애, 사회·경제적 상황 등 취약성을 악용하는 AI 시스템
 \item \textbf{사회적 점수화(Social Scoring)}: 공공기관이 사회적 행동에 기반하여 개인을 평가·분류하는 시스템
 \item \textbf{개인 범죄 예측(Individual Criminal Prediction)}: 프로파일링에만 기반한 개인 범죄 리스크 평가
 \item \textbf{무차별적 안면 인식 데이터베이스}: 인터넷·CCTV 영상의 무차별적 수집을 통한 안면 인식 데이터베이스 구축
 \item \textbf{감정 인식(Emotion Recognition)}: 직장 및 교육 기관에서의 감정 인식 시스템 (안전 또는 의료 목적 제외)
 \item \textbf{생체인식 분류(Biometric Categorisation)}: 인종, 정치 성향, 종교, 성적 지향 등 민감 속성에 기반한 분류
 \item \textbf{실시간 원격 생체 인식(Real-time Remote Biometric Identification)}: 법 집행 목적의 공공장소 실시간 원격 생체 인식 (엄격한 예외 사유 존재)
\end{enumerate}

\paragraph{2단계: 고위험 AI 시스템(High-Risk)}
고위험 AI 시스템은 가장 광범위하고 상세한 규제 의무가 부과되는 범주이다. 2026년 8월 2일부터 본격 적용된다. Annex III에 정의된 8개 영역은 다음과 같다:

\begin{enumerate}
 \item \textbf{생체인식 및 생체 기반 분류}: 원격 생체 인식 시스템, 감정 인식 시스템 (금지 범주에 해당하지 않는 것)
 \item \textbf{핵심 인프라 관리 및 운영}: 도로 교통, 수도·가스·난방·전기 공급의 안전 구성요소
 \item \textbf{교육 및 직업 훈련}: 입학 결정, 학생 평가, 학습 수준 평가에 사용되는 AI
 \item \textbf{고용, 인력 관리, 자영업 접근}: 채용 및 선발, 근무 조건 결정, 승진 및 해고 결정에 사용되는 AI
 \item \textbf{필수 민간 서비스 접근}: 신용 평가, 생명·건강 보험 리스크 평가, 긴급 서비스 우선순위 결정
 \item \textbf{법 집행}: 범죄 분석, 리스크 평가, 거짓말 탐지, 증거 신뢰성 평가에 사용되는 AI
 \item \textbf{이민, 망명, 국경 통제}: 리스크 평가, 비자·거주 허가 신청 심사, 이민 관련 민원 처리에 사용되는 AI
 \item \textbf{사법 및 민주적 과정}: 법적 해석, 사실 관계 적용, 대체적 분쟁 해결에 사용되는 AI
\end{enumerate}

고위험 AI 시스템에 부과되는 주요 의무는 다음과 같다:
\begin{itemize}
 \item 리스크관리 시스템(Risk Management System) 구축
 \item 데이터 거버넌스(Data Governance) 체계 수립
 \item 기술 문서(Technical Documentation) 작성
 \item 기록 보관(Record-keeping) 및 로깅(Logging) 체계 구축
 \item 사용자에 대한 투명성 확보 및 정보 제공
 \item 인간 감독(Human Oversight) 메커니즘 구현
 \item 정확성, 견고성, 사이버보안 확보
 \item 적합성 평가(Conformity Assessment) 통과\cite{mokander2022conformity}
\end{itemize}

\paragraph{3단계: 제한적 리스크 AI 시스템(Limited Risk)}
투명성 의무만 부과되는 범주이다. 챗봇, 감정 인식 시스템, 딥페이크(Deepfake) 생성 시스템 등이 해당한다. 사용자에게 AI와 상호작용하고 있음을 고지하고, AI 생성 콘텐츠임을 표시해야 한다.

\paragraph{4단계: 최소 리스크 AI 시스템(Minimal Risk)}
대부분의 AI 시스템이 해당하며, 추가적인 법적 의무는 부과되지 않는다. 스팸 필터, AI 기반 게임 등이 이에 해당한다. 다만 자발적 행동강령(Code of Conduct) 준수를 권장한다.

\subsubsection{적합성 평가(Conformity Assessment) 절차}

고위험 AI 시스템은 시장 출시 전 적합성 평가를 통과해야 한다. 평가 방식은 두 가지이다:

\begin{itemize}
 \item \textbf{자체 적합성 평가(Self-Assessment)}: 대부분의 고위험 AI 시스템에 적용. 공급자가 자체적으로 기술 문서를 작성하고, EU 적합성 선언(Declaration of Conformity)을 발행하며, CE 마킹을 부착한다.
 \item \textbf{제3자 적합성 평가(Third-Party Assessment)}: 생체 인식 관련 고위험 AI 시스템에 적용. 지정된 적합성 평가 기관(Notified Body)이 독립적으로 평가를 수행한다.
\end{itemize}

\subsubsection{범용 AI(GPAI) 규정 상세}

2025년 8월 2일부터 적용되는 범용 AI 모델(General-Purpose AI Model) 규정은 기반 모델(Foundation Model)과 대규모 언어 모델(LLM)에 대한 의무를 규정한다.

\paragraph{모든 GPAI 모델 공급자의 의무}
\begin{itemize}
 \item 기술 문서 작성 및 유지
 \item EU AI Office에 정보 제공
 \item 저작권 관련 정책 수립 (학습 데이터의 저작권법 준수)
 \item 학습 데이터에 대한 충분히 상세한 요약(Summary) 공개
\end{itemize}

\paragraph{시스템적 리스크(Systemic Risk) 모델 기준}
GPAI 모델 중 시스템적 리스크를 보유한 모델에는 추가 의무가 부과된다. 시스템적 리스크 모델의 판단 기준은 다음과 같다:
\begin{itemize}
 \item \textbf{연산량 기준}: 학습에 $10^{25}$ FLOPs(부동소수점 연산) 이상이 사용된 모델
 \item \textbf{EU 집행위원회 지정}: 연산량 기준에 해당하지 않더라도, 높은 영향력(High Impact Capability)을 보유한 것으로 판단되는 모델을 지정할 수 있음
\end{itemize}

시스템적 리스크 모델 공급자의 추가 의무:
\begin{itemize}
 \item 모델 평가(Model Evaluation) 수행: 표준화된 프로토콜에 따른 체계적 평가
 \item 시스템적 리스크 평가 및 완화 조치 이행
 \item 심각한 인시던트(Serious Incident) 추적 및 보고 체계 구축
 \item 적절한 사이버보안 보호 수준 확보
 \item AI Office에 에너지 소비 정보 보고
\end{itemize}

\subsubsection{제재 및 과징금 체계}

EU AI Act는 위반 유형에 따라 차등적 과징금을 부과한다:

\begin{center}
\renewcommand{\arraystretch}{1.3}
\small
\begin{tabularx}{\textwidth}{Xp{3cm}p{3cm}}
\toprule
\textbf{위반 유형} & \textbf{과징금 상한} & \textbf{중소기업 상한} \\
\midrule
금지 AI 시스템 운영 & 3,500만 유로 또는 전 세계 매출의 7\% & 750만 유로 또는 매출의 1\% \\
고위험 AI 의무 위반 & 1,500만 유로 또는 전 세계 매출의 3\% & 750만 유로 또는 매출의 1\% \\
부정확한 정보 제공 & 750만 유로 또는 전 세계 매출의 1\% & 750만 유로 또는 매출의 1\% \\
GPAI 의무 위반 & 1,500만 유로 또는 전 세계 매출의 3\% & 해당 비례 적용 \\
\bottomrule
\end{tabularx}
\end{center}

\subsubsection{EU AI Act 시행 타임라인}

\begin{table}[htbp]
\centering
\caption{EU AI Act 시행 타임라인 (2026년 5월 잠정 합의 반영)}
\label{tab:eu_ai_act_timeline_global}
\begin{tabularx}{\textwidth}{p{2.5cm}p{3.2cm}X}
\toprule
\textbf{시점} & \textbf{적용 영역} & \textbf{주요 내용} \\
\midrule
2024.8.1 & 발효 & EU AI Act 공식 발효 \\
2025.2.2 & 금지 AI & 금지 AI 시스템 규정 시행, AI 리터러시 의무 \\
2025.8.2 & GPAI & 범용 AI 모델 규정 적용, 거버넌스 구조 확립 \\
2026.8.2 & \textit{(원래 기한)} & \textit{고위험 AI Annex III 기한 — Digital Omnibus로 연기 추진 중} \\
2026.12.2 & AI 워터마킹 & AI 생성 콘텐츠 워터마킹 의무 적용 (잠정 합의) \\
2027.12.2 & 고위험 AI (Annex III) & 독립형 고위험 AI 시스템 의무 — 2026.5.7 EU 이사회·의회 잠정 합의로 연기 \\
2028.8.2 & 고위험 AI (Annex I) & 의료기기·기계류 등 기존 EU 법률 규제 제품 내 고위험 AI \\
\bottomrule
\end{tabularx}
\end{table}

\begin{tcolorbox}[colback=orange!5, colframe=orange!60, title={2026년 5월 최신 동향: EU AI Act Digital Omnibus}]
2026년 5월 7일 EU 이사회와 유럽의회는 Digital Omnibus 이니셔티브 내 AI Act 조항에 대한 \textbf{잠정 합의(provisional agreement)}에 도달했다. 핵심 변경 사항은 다음과 같다:
\begin{itemize}
    \item \textbf{Annex III 고위험 AI 기한 연기}: 2026년 8월 → \textbf{2027년 12월 2일}
    \item \textbf{Annex I(의료기기·기계류) 기한 연기}: 2027년 8월 → \textbf{2028년 8월 2일}
    \item \textbf{NCII/CSAM 신규 금지}: AI 기반 비동의 성적 이미지(NCII) 및 아동 성착취물(CSAM) 생성 AI 시스템 금지 조항 신설
    \item \textbf{AI 워터마킹 의무}: AI 생성 콘텐츠 식별 의무 2026년 12월 2일 적용
\end{itemize}
단, 잠정 합의는 공식 입법 채택 절차가 남아 있으며, 채택 전까지는 원래 일정(2026년 8월)이 법적으로 유효하다. EU 시장에 AI 시스템을 제공하는 한국 기업은 진행 상황을 면밀히 추적해야 한다.
\end{tcolorbox}

\subsubsection{한국 기업의 EU AI Act 준수 전략}

EU 시장에 AI 제품·서비스를 제공하는 한국 기업은 역외 적용(Extraterritorial Application) 조항에 따라 EU AI Act를 준수해야 한다. 핵심 준수 전략은 다음과 같다:

\begin{enumerate}
 \item \textbf{AI 시스템 인벤토리 및 리스크 분류}: EU 시장에서 운영되는 모든 AI 시스템을 목록화하고, 4단계 리스크 분류에 따라 분류한다.
 \item \textbf{EU 대리인(Authorised Representative) 지정}: EU 내에 법적 대리인을 지정하여 규제 당국과의 소통 채널을 확보한다.
 \item \textbf{기술 문서 체계 구축}: 고위험 AI 시스템에 대한 기술 문서(설계 사양, 학습 데이터 설명, 테스트 결과, 리스크 분석)를 체계적으로 작성·유지한다.
 \item \textbf{적합성 평가 준비}: 조화 표준(Harmonised Standards)에 따른 적합성 평가를 준비하고, 필요시 적합성 평가 기관(Notified Body)과 사전 협의한다.
 \item \textbf{GPAI 의무 이행}: 범용 AI 모델을 제공하는 경우, 기술 문서 작성, 저작권 정책 수립, 학습 데이터 요약 공개 등의 의무를 이행한다.
\end{enumerate}

%----------------------------------------------------------
\subsection{미국: 연방 주도 자율 규제로의 전환}
%----------------------------------------------------------
2025년 트럼프 행정부 출범으로 규제 기조가 급변하였다.
\begin{itemize}
 \item \textbf{EO 14179 (2025.1)}: 바이든 행정부의 EO 14110 폐지, 혁신 촉진으로 선회. AI 안전 보고 의무 등 기존 규제 요건을 철폐하고, 미국 AI 산업의 글로벌 리더십 유지를 최우선 목표로 설정.
 \item \textbf{EO 14199 (2025.12)}: ``연방 프리엠션 전략'' 추진. DOJ(법무부) 소송 TF 설립, 연방 자금 연계를 통해 주(州) 단위 규제(캘리포니아 SB 1047, 콜로라도 SB 24-205 등)를 무력화하고 연방 차원의 통일된(완화된) 기준 추진.
 \item \textbf{주(州) 규제 동향}: 연방 프리엠션에도 불구하고, 캘리포니아, 뉴욕, 일리노이 등 주요 주에서는 독자적 AI 규제를 추진 중이며, 법적 분쟁이 진행 중이다.
\end{itemize}

한국 기업의 대응 전략으로는, 연방-주 규제 간 불확실성이 높은 만큼 보수적 접근이 필요하며, 가장 엄격한 주 규제 기준을 기본 준수선으로 설정하는 것이 리스크 관리 측면에서 유리하다.

2026년 3월 백악관이 \textbf{국가 AI 정책 프레임워크(National Policy Framework for AI)}를 발표하고, 연방법으로 주법을 선점(preemption)하도록 의회에 권고하면서 연방-주 충돌이 본격화되었다. 동시에 Colorado AI Act, California TFAIA, Texas RAIGA 등 주 단위 AI 입법이 활발히 추진되고 있어, 미국 진출 기업은 연방 수준과 진출 주 별도 규제 모두를 동시에 추적해야 하는 복잡한 환경에 놓여 있다.

%----------------------------------------------------------
\subsection{중국: 알고리즘 거버넌스와 콘텐츠 라벨링 강화}
%----------------------------------------------------------
\begin{itemize}
 \item \textbf{AI 콘텐츠 라벨링 (2025.9.1 시행)}: 명시적 라벨(가시적 표시)과 암시적 라벨(워터마크 등) 의무화. 위반 시 강력한 단속. AI 서비스 제공자뿐만 아니라 콘텐츠 배포 플랫폼에도 라벨링 의무를 부과한다.
 \item \textbf{기존 규제}: 알고리즘 추천 관리 규정(2022), 딥 합성 관리 규정(2023), 생성형 AI 관리 잠정 방법(2023) 등 분야별 촘촘한 규제망을 운영하고 있다.
 \item \textbf{개정 사이버보안법 (2026.1.1 발효)}: AI 거버넌스 조항이 중국 기본법 수준 입법에 최초로 명문화되었다.
 \item \textbf{의인화 AI 상호작용 서비스 규제 (2026.7.15 시행 예정)}: 국가사이버공간관리국(CAC) 등 5개 기관이 공동 발표한 AI 동반자·에이전트 서비스 규제로, 미성년자 대상 AI 동반자 서비스를 금지하고 의인화 AI의 감정 조작을 규율한다.
 \item \textbf{특징}: 콘텐츠 통제와 사회 안정을 우선시하며, 서비스 출시 전 ``알고리즘 등록(Algorithm Filing)'' 의무를 부과하는 등 사전 규제 성격이 강하다.\cite{roberts2021chinese}
\end{itemize}

%----------------------------------------------------------
\subsection{일본: AI 진흥법 제정과 소프트 로(Soft Law) 모델}
%----------------------------------------------------------
일본은 가이드라인 중심의 소프트 로 접근을 선호하다가, 2025년 최초의 AI 전문 법률을 제정하였다.\cite{jobin2019global}
\begin{itemize}
 \item \textbf{AI 진흥법 (2025.6.4 발효)}: 일본 최초의 AI 전문 법률. 단, EU AI Act와 달리 벌금 없는 ``권고→정보 요청→공개(name and shame)'' 방식이며, 혁신 촉진을 우선 목표로 설계되었다. AI의 연구개발 및 활용 촉진을 국가 과제로 명문화한 것이 핵심이다.
 \item \textbf{APPI 개정안 (2026.4 국회 제출)}: AI 관련 데이터 이용 일부 완화를 포함한 개인정보보호법 개정안이 정부 승인을 받아 국회에 제출되었다.
 \item \textbf{AI 사업자 가이드라인 (2024.4)}: 인간 중심, 안전성, 공정성, 프라이버시 보호, 보안, 투명성, 책임성, 교육·리터러시, 공정 경쟁, 혁신의 10대 원칙을 제시한다.
 \item \textbf{히로시마 AI 프로세스}: G7 의장국으로서 국제 AI 거버넌스 논의를 주도하며, 자발적 행동강령(Code of Conduct) 기반의 국제 협력을 추구한다.
\end{itemize}

%----------------------------------------------------------
\subsection{싱가포르: 실용주의적 거버넌스 프레임워크}
%----------------------------------------------------------
\begin{itemize}
 \item \textbf{Model AI 거버넌스 프레임워크}: 실용적 지침을 제공하며, 설명가능성, 투명성, 인간 감독 등을 핵심 원칙으로 제시한다.
 \item \textbf{AI Verify}: 거버넌스 테스트 툴킷을 오픈소스로 공개하여\cite{singapore2020model}, 기업이 자체적으로 AI 시스템의 거버넌스 수준을 검증할 수 있도록 한다. 공정성, 견고성, 설명가능성 등 11개 영역에 대한 자동화된 테스트를 제공한다.
 \item \textbf{GenAI 거버넌스 프레임워크}: 생성형 AI에 특화된 거버넌스 프레임워크를 추가로 발표하여, 환각, 저작권, 콘텐츠 안전 등의 이슈에 대한 가이드를 제공한다.
\end{itemize}

%==========================================================
\section{국내 AI 규제 환경}
%==========================================================

%----------------------------------------------------------
\subsection{AI 기본법 심층 분석}
%----------------------------------------------------------
2024년 12월 국회를 통과하여 2026년 1월 22일 시행되는 ``AI 기본법(인공지능 산업 육성 및 신뢰 기반 조성 등에 관한 법률)''은 한국 AI 거버넌스의 토대이다.\cite{cset2025korea_ai_act}

\subsubsection{입법 배경 및 기본 방향}
AI 기본법은 ``진흥 우선, 규제 최소화'' 원칙 하에 제정되었다. 1년의 유예 기간을 두어 기업의 준비 부담을 완화하면서도, 고위험 AI에 대한 최소한의 안전장치를 마련하였다. EU AI Act와 비교하면 규제 강도는 낮지만, 아시아 국가 중에서는 선도적인 포괄적 AI 법률이다.

\subsubsection{주요 조항 상세 분석}

\paragraph{고위험 AI 분류 기준 (11개 섹터)}
AI 기본법은 다음 11개 분야에서 사용되는 AI 시스템을 고위험 AI로 분류한다:

\begin{enumerate}
 \item \textbf{생명·신체의 안전}: 의료 진단, 치료 지원, 의료기기에 사용되는 AI
 \item \textbf{기본권 침해}: 개인의 기본적 권리에 중대한 영향을 미치는 AI
 \item \textbf{사법·수사}: 범죄 수사, 재판 지원에 사용되는 AI
 \item \textbf{출입국 관리}: 비자 심사, 입국 심사 지원에 사용되는 AI
 \item \textbf{교육}: 입학 심사, 학업 성취 평가에 사용되는 AI
 \item \textbf{고용·노동}: 채용, 인사 평가, 해고 결정에 사용되는 AI
 \item \textbf{금융}: 신용 평가, 보험 심사, 투자 자문에 사용되는 AI
 \item \textbf{의료}: 질병 진단, 처방 추천, 건강 리스크 예측에 사용되는 AI
 \item \textbf{복지}: 사회 복지 수급 자격 판정에 사용되는 AI
 \item \textbf{안전}: 핵심 인프라(전력, 교통, 통신) 관리에 사용되는 AI
 \item \textbf{대규모 모델}: 학습에 $10^{26}$ FLOPs 이상이 사용된 초거대 AI 모델
\end{enumerate}

\paragraph{영향평가 의무}
고위험 AI 시스템의 개발자 및 운영자는 다음 사항을 포함하는 영향평가를 수행해야 한다:
\begin{itemize}
 \item AI 시스템의 목적, 기능, 한계에 대한 설명
 \item 잠재적 리스크 및 부정적 영향 분석
 \item 리스크 완화 조치 및 안전 확보 방안
 \item 데이터의 적정성 및 편향 검토 결과
 \item 인간 감독 체계 및 이의신청 절차
\end{itemize}

\paragraph{설명 의무}
고위험 AI 시스템에 의해 영향을 받는 개인은 해당 결정에 대한 설명을 요청할 권리가 있다. AI 운영자는 다음을 설명해야 한다:
\begin{itemize}
 \item AI 시스템이 의사결정에 관여하였다는 사실
 \item 해당 의사결정의 주요 판단 근거
 \item 이의신청 및 구제를 받을 수 있는 절차
\end{itemize}

\paragraph{생성형 AI 표시 의무}
생성형 AI로 생성된 콘텐츠(텍스트, 이미지, 영상, 음성)에 대해 다음의 표시 의무가 부과된다:
\begin{itemize}
 \item AI에 의해 생성된 콘텐츠임을 명확히 표시
 \item 워터마크(Watermark) 등 기술적 조치를 통한 식별 수단 제공
 \item 딥페이크 콘텐츠의 경우, 실제가 아닌 합성 콘텐츠임을 고지
\end{itemize}

\paragraph{벌칙 및 과징금 체계}
\begin{itemize}
 \item 고위험 AI 의무 위반: 3억원 이하의 과징금
 \item 생성형 AI 표시 의무 위반: 1억원 이하의 과태료
 \item 영향평가 미수행: 2억원 이하의 과징금
 \item 시정명령 불이행: 5억원 이하의 과징금
\end{itemize}

\subsubsection{AI 기본법 시행 타임라인}

\begin{table}[htbp]
\centering
\caption{한국 AI 기본법 시행 타임라인}
\label{tab:korea_ai_basic_timeline}
\begin{tabularx}{\textwidth}{p{2.5cm}X}
\toprule
\textbf{시점} & \textbf{주요 내용} \\
\midrule
2024.12 & 국회 본회의 통과 \\
2025.1 & 대통령 공포 \\
2025.1$\sim$12 & 시행령·시행규칙 제정, 고위험 AI 세부 기준 마련, 기업 준비 기간 \\
2026.1.22 & AI 기본법 시행 \\
2026.상반기 & 고위험 AI 신고 의무 개시, 영향평가 체계 본격 가동 \\
2026.하반기 & 생성형 AI 표시 의무 시행, 과징금 부과 체계 가동 \\
\bottomrule
\end{tabularx}
\end{table}

%----------------------------------------------------------
\subsection{기타 주요 규제}
%----------------------------------------------------------
\begin{itemize}
 \item \textbf{개인정보보호법}: 자동화된 의사결정에 대한 설명·거부권 도입. 2025년 3월 시행된 개정법은 완전 자동화된 의사결정에 대한 거부권과 인적 개입 요구권을 규정한다.
 \item \textbf{금융 분야}: 금융위원회 ``AI 활용 가이드라인''\cite{fsc2024ai}은 금융 AI의 공정성, 투명성, 안전성에 대한 세부 기준을 제시한다. 신용평가 AI에 대한 특별 감독 체계를 운영한다.
 \item \textbf{의료 분야}: 식약처 ``SaMD(Software as a Medical Device) 가이드라인''은 의료 AI의 인허가 절차, 임상 검증 요건, 시판 후 감시 체계를 규정한다.
 \item \textbf{공공 분야}: 행정안전부 ``공공부문 AI 활용 가이드라인''\cite{nia2024quality}은 공공기관의 AI 도입 시 투명성, 참여, 책임성 원칙을 제시한다.
\end{itemize}

%==========================================================
\section{규제 수렴과 상호운용성}
%==========================================================
글로벌 규제는 투명성, 책임성 등 원칙 수준에서는 수렴하나\cite{hagendorff2020ethics, whittlestone2019role}, 구체적 실행(EU의 리스크 기반 규제 vs 미국의 자율 규제)에서는 분화되고 있다. 특히 트럼프 행정부의 연방 프리엠션 전략은 미국-EU 간 규제 철학 충돌을 심화시켰다.

다국적 기업은 통합 거버넌스 전략이 필요하다. 핵심 원칙은 최고 수준(EU AI Act)을 따르되, 세부 절차는 관할권별 모듈화를 통해 유연하게 대응해야 한다.

%----------------------------------------------------------
\subsection{규제 격차 분석(Gap Analysis) 방법론}
%----------------------------------------------------------

규제 격차 분석은 조직의 현재 AI 거버넌스 수준과 목표 규제 요건 간의 차이를 체계적으로 파악하는 프로세스이다. 다음 5단계로 수행한다:

\begin{enumerate}
 \item \textbf{규제 요건 매핑}: 적용 대상 규제(EU AI Act, AI 기본법 등)의 구체적 요건을 항목별로 분해한다.
 \item \textbf{현행 수준 평가}: 각 요건 항목에 대한 조직의 현재 준수 수준을 ``미준수/부분 준수/실질적 준수/완전 준수''의 4단계로 평가한다.
 \item \textbf{격차 식별 및 우선순위화}: 미준수 또는 부분 준수 항목을 식별하고, 리스크 영향도와 시행 시점을 고려하여 우선순위를 매긴다.
 \item \textbf{대응 계획 수립}: 격차 해소를 위한 구체적 조치, 담당자, 일정, 필요 자원을 정의한다.
 \item \textbf{이행 및 추적}: 대응 계획을 이행하고, 정기적으로 진척 상황을 추적한다.
\end{enumerate}

%----------------------------------------------------------
\subsection{다국적 기업의 통합 거버넌스 전략}
%----------------------------------------------------------

다국적 기업이 복수의 관할권에서 AI 시스템을 운영할 때, 각 관할권의 규제를 개별적으로 대응하는 것은 비효율적이다. 통합 거버넌스 전략의 핵심 원칙은 다음과 같다:

\begin{enumerate}
 \item \textbf{최고 기준 원칙(Highest Standard Principle)}: 적용 대상 규제 중 가장 엄격한 기준(현재로서는 EU AI Act)을 글로벌 기본선(Baseline)으로 설정한다.\cite{bradford2020brussels} 이를 통해 어떤 관할권에서도 준수를 보장한다.
 \item \textbf{모듈화된 준수 체계}: 글로벌 기본선 위에 관할권별 특수 요건을 모듈로 추가한다. 예를 들어, 중국 시장에는 콘텐츠 라벨링 모듈을, 한국 시장에는 AI 기본법 고위험 AI 신고 모듈을 추가한다.
 \item \textbf{중앙 집중식 거버넌스 + 지역 자율성}: 글로벌 AI 거버넌스 정책과 리스크관리 프레임워크는 본사에서 수립하되, 지역별 이행은 현지 법규 전문가의 자문을 받아 유연하게 운영한다.
 \item \textbf{규제 변화 추적 체계}: 주요 관할권의 규제 동향을 실시간으로 모니터링하고, 변화가 감지되면 영향 분석을 수행하는 체계를 구축한다.
\end{enumerate}

%==========================================================
\section{규제 대응 실무 도구}
%==========================================================

%----------------------------------------------------------
\subsection{글로벌 AI 규제 비교 매트릭스}
%----------------------------------------------------------

\begin{practicaltool}{글로벌 AI 규제 비교 매트릭스 (2026년 1월 기준)}
본 매트릭스는 주요 관할권의 규제 성격을 비교하여 대응 전략 수립을 지원한다.

\begin{center}
\renewcommand{\arraystretch}{1.3}
\small
\begin{tabularx}{\textwidth}{p{1.8cm}p{2.2cm}p{2.2cm}p{2.2cm}p{2.2cm}p{2.2cm}}
\toprule
\textbf{구분} & \textbf{EU} & \textbf{미국} & \textbf{중국} & \textbf{한국} & \textbf{싱가포르} \\
\midrule
규제 성격 & 구속력 있는 법률 & 자발적/ 프리엠션 추진 & 구속력 있는 규정 & 법률 (26.1 시행) & 자발적 가이드 \\
\midrule
접근법 & 리스크 기반 & 혁신 촉진 & 분야별 규제 & 리스크 기반 & 실용주의 \\
\midrule
고위험 정의 & Annex III 8개 영역 & 기준 없음 & 명시적 정의 없음 & 11개 섹터 + $10^{26}$ FLOPs & 가이드 기반 \\
\midrule
금지 AI & 8개 유형 & 없음 & 일부 활용 금지 & 명시적 없음 & 없음 \\
\midrule
라벨링 & 의무 (26.8 예정) & 없음 & 의무 (25.9 시행) & 의무 (생성형) & 자발적 \\
\midrule
GPAI 규정 & $10^{25}$ FLOPs 기준 & 없음 & 등록 의무 & $10^{26}$ FLOPs & 없음 \\
\midrule
과징금 & 매출 7\% & 해당 없음 & 행정 처분 & 5억원 이하 & 해당 없음 \\
\midrule
역외 적용 & 있음 & 없음 & 서비스 기준 & 제한적 & 없음 \\
\bottomrule
\end{tabularx}
\end{center}
\vspace{0.2cm}
\textbf{활용 방법}:
\begin{enumerate}
 \item 적용 관할권 식별 및 AI 시스템 인벤토리 매핑
 \item 시스템별 리스크 분류 및 규제 의무 격차(Gap) 분석
 \item 미국 사업 시 연방 프리엠션(주법 무력화) 리스크 별도 평가
 \item 중국 사업 시 알고리즘 등록 및 콘텐츠 라벨링 의무 확인
 \item 분기별 업데이트 및 규제 변화 추적
\end{enumerate}
\end{practicaltool}

%----------------------------------------------------------
\subsection{프레임워크 선택 가이드}
%----------------------------------------------------------

조직의 상황에 따라 우선적으로 적용할 프레임워크를 선택하는 가이드를 제공한다.

\begin{practicaltool}{AI 거버넌스 프레임워크 선택 가이드}

\begin{center}
\renewcommand{\arraystretch}{1.3}
\small
\begin{tabularx}{\textwidth}{p{3cm}p{3.5cm}X}
\toprule
\textbf{조직 유형} & \textbf{우선 프레임워크} & \textbf{권고 사항} \\
\midrule
EU 시장 진출 기업 & EU AI Act + ISO 42001 & 적합성 평가 준비를 최우선으로 하고, ISO 42001 인증으로 준수 역량을 입증한다 \\
\midrule
미국 시장 중심 기업 & NIST AI RMF + ISO 42001 & NIST AI RMF로 리스크관리 체계를 수립하되, 주별 규제 동향을 면밀히 추적한다 \\
\midrule
국내 사업 중심 기업 & AI 기본법 + NIST AI RMF & AI 기본법 준수를 기본으로 하고, NIST AI RMF로 리스크관리 체계를 보강한다 \\
\midrule
글로벌 다국적 기업 & EU AI Act 기반 통합 + ISO 42001 & EU AI Act를 글로벌 기본선으로 설정하고, 관할권별 모듈을 추가한다 \\
\midrule
AI 모델 공급자 & EU AI Act GPAI + NIST GenAI Profile & 범용 AI 모델 의무와 생성형 AI 특화 리스크에 집중한다 \\
\midrule
스타트업·중소기업 & OECD 원칙 + AI 기본법 & 원칙 수준의 체계를 먼저 수립하고, 사업 확장에 따라 단계적으로 강화한다 \\
\bottomrule
\end{tabularx}
\end{center}
\end{practicaltool}

%----------------------------------------------------------
\subsection{규제 매핑 자동화 도구}
%----------------------------------------------------------

\begin{practicaltool}{규제 매핑 자동화 도구}
다음 Python 코드는 다중 규제 환경에서 AI 시스템별 적용 규제를 자동으로 매핑하고, 격차를 분석하는 도구이다.

\begin{lstlisting}[language=Python]
"""Multi-Jurisdiction AI Regulation Mapping Tool"""
from dataclasses import dataclass, field
from typing import Dict, List, Set, Optional
from enum import Enum
import json

class RiskLevel(Enum):
    MINIMAL = "minimal"
    LIMITED = "limited"
    HIGH = "high"
    UNACCEPTABLE = "unacceptable"

class Jurisdiction(Enum):
    EU = "EU"
    US = "US"
    CHINA = "China"
    KOREA = "Korea"
    SINGAPORE = "Singapore"
    JAPAN = "Japan"

class ComplianceStatus(Enum):
    NOT_ASSESSED = "not_assessed"
    NON_COMPLIANT = "non_compliant"
    PARTIAL = "partial"
    SUBSTANTIAL = "substantial"
    COMPLIANT = "compliant"

@dataclass
class RegulatoryRequirement:
    id: str
    jurisdiction: Jurisdiction
    description: str
    risk_level: RiskLevel
    deadline: str
    mandatory: bool = True

@dataclass
class AISystem:
    name: str
    description: str
    domain: str
    jurisdictions: Set[Jurisdiction] = field(
        default_factory=set
    )
    risk_classification: Dict[Jurisdiction, RiskLevel] = (
        field(default_factory=dict)
    )

class RegulatoryMapper:
    """Maps AI systems to applicable regulations."""

    def __init__(self):
        self.requirements: List[RegulatoryRequirement] = []
        self.ai_systems: List[AISystem] = []
        self._load_requirements()

    def _load_requirements(self):
        """Load regulatory requirements database."""
        # EU AI Act requirements
        eu_reqs = [
            RegulatoryRequirement(
                "EU-HIGH-001", Jurisdiction.EU,
                "Risk management system",
                RiskLevel.HIGH, "2026-08-02",
            ),
            RegulatoryRequirement(
                "EU-HIGH-002", Jurisdiction.EU,
                "Data governance measures",
                RiskLevel.HIGH, "2026-08-02",
            ),
            RegulatoryRequirement(
                "EU-HIGH-003", Jurisdiction.EU,
                "Technical documentation",
                RiskLevel.HIGH, "2026-08-02",
            ),
            RegulatoryRequirement(
                "EU-HIGH-004", Jurisdiction.EU,
                "Human oversight mechanism",
                RiskLevel.HIGH, "2026-08-02",
            ),
            RegulatoryRequirement(
                "EU-HIGH-005", Jurisdiction.EU,
                "Conformity assessment",
                RiskLevel.HIGH, "2026-08-02",
            ),
            RegulatoryRequirement(
                "EU-GPAI-001", Jurisdiction.EU,
                "Technical documentation for GPAI",
                RiskLevel.LIMITED, "2025-08-02",
            ),
            RegulatoryRequirement(
                "EU-GPAI-002", Jurisdiction.EU,
                "Copyright compliance policy",
                RiskLevel.LIMITED, "2025-08-02",
            ),
            RegulatoryRequirement(
                "EU-TRANS-001", Jurisdiction.EU,
                "AI interaction transparency",
                RiskLevel.LIMITED, "2026-08-02",
            ),
        ]
        self.requirements.extend(eu_reqs)

        # Korea AI Basic Act requirements
        kr_reqs = [
            RegulatoryRequirement(
                "KR-HIGH-001", Jurisdiction.KOREA,
                "Impact assessment for high-risk AI",
                RiskLevel.HIGH, "2026-07-01",
            ),
            RegulatoryRequirement(
                "KR-HIGH-002", Jurisdiction.KOREA,
                "Explanation obligation",
                RiskLevel.HIGH, "2026-07-01",
            ),
            RegulatoryRequirement(
                "KR-HIGH-003", Jurisdiction.KOREA,
                "Safety measures for high-risk AI",
                RiskLevel.HIGH, "2026-07-01",
            ),
            RegulatoryRequirement(
                "KR-GEN-001", Jurisdiction.KOREA,
                "Generative AI content labeling",
                RiskLevel.LIMITED, "2026-07-01",
            ),
        ]
        self.requirements.extend(kr_reqs)

        # China requirements
        cn_reqs = [
            RegulatoryRequirement(
                "CN-LABEL-001", Jurisdiction.CHINA,
                "AI content explicit labeling",
                RiskLevel.LIMITED, "2025-09-01",
            ),
            RegulatoryRequirement(
                "CN-LABEL-002", Jurisdiction.CHINA,
                "AI content implicit labeling (watermark)",
                RiskLevel.LIMITED, "2025-09-01",
            ),
            RegulatoryRequirement(
                "CN-ALGO-001", Jurisdiction.CHINA,
                "Algorithm registration filing",
                RiskLevel.LIMITED, "2022-03-01",
            ),
        ]
        self.requirements.extend(cn_reqs)

    def register_system(self, system: AISystem):
        """Register an AI system for regulatory mapping."""
        self.ai_systems.append(system)

    def get_applicable_requirements(
        self, system: AISystem
    ) -> List[RegulatoryRequirement]:
        """Get all requirements applicable to a system."""
        applicable = []
        for req in self.requirements:
            if req.jurisdiction not in system.jurisdictions:
                continue
            sys_risk = system.risk_classification.get(
                req.jurisdiction, RiskLevel.MINIMAL
            )
            risk_order = [
                RiskLevel.MINIMAL, RiskLevel.LIMITED,
                RiskLevel.HIGH, RiskLevel.UNACCEPTABLE,
            ]
            if (risk_order.index(sys_risk)
                    >= risk_order.index(req.risk_level)):
                applicable.append(req)
        return applicable

    def gap_analysis(
        self,
        system: AISystem,
        current_status: Dict[str, ComplianceStatus],
    ) -> dict:
        """Perform gap analysis for a system."""
        applicable = self.get_applicable_requirements(system)
        gaps = []
        compliant_count = 0

        for req in applicable:
            status = current_status.get(
                req.id, ComplianceStatus.NOT_ASSESSED
            )
            if status in (
                ComplianceStatus.COMPLIANT,
                ComplianceStatus.SUBSTANTIAL,
            ):
                compliant_count += 1
            else:
                gaps.append({
                    "requirement_id": req.id,
                    "jurisdiction": req.jurisdiction.value,
                    "description": req.description,
                    "deadline": req.deadline,
                    "current_status": status.value,
                    "mandatory": req.mandatory,
                })

        total = len(applicable)
        return {
            "system": system.name,
            "total_requirements": total,
            "compliant": compliant_count,
            "compliance_rate": (
                round(compliant_count / total * 100, 1)
                if total > 0 else 0
            ),
            "gaps": sorted(
                gaps, key=lambda x: x["deadline"]
            ),
        }


# Usage example
if __name__ == "__main__":
    mapper = RegulatoryMapper()

    # Register a credit scoring AI system
    credit_ai = AISystem(
        name="AI Credit Scoring System",
        description="ML-based credit evaluation",
        domain="finance",
        jurisdictions={
            Jurisdiction.EU, Jurisdiction.KOREA,
        },
        risk_classification={
            Jurisdiction.EU: RiskLevel.HIGH,
            Jurisdiction.KOREA: RiskLevel.HIGH,
        },
    )
    mapper.register_system(credit_ai)

    # Current compliance status
    status = {
        "EU-HIGH-001": ComplianceStatus.SUBSTANTIAL,
        "EU-HIGH-002": ComplianceStatus.PARTIAL,
        "EU-HIGH-003": ComplianceStatus.PARTIAL,
        "KR-HIGH-001": ComplianceStatus.NOT_ASSESSED,
    }

    result = mapper.gap_analysis(credit_ai, status)
    print(json.dumps(result, indent=2, ensure_ascii=False))
    print(f"\nCompliance Rate: {result['compliance_rate']}%")
    print(f"Gaps to address: {len(result['gaps'])}")
\end{lstlisting}
\end{practicaltool}

%----------------------------------------------------------
\subsection{규제 변화 추적 체계}
%----------------------------------------------------------

글로벌 AI 규제 환경은 빠르게 변화하고 있으며, 이를 체계적으로 추적하고 대응하는 체계가 필수적이다. 효과적인 규제 변화 추적 체계는 다음 요소로 구성된다:

\begin{enumerate}
 \item \textbf{모니터링 대상 정의}: 조직에 영향을 미치는 관할권의 법률, 규정, 가이드라인, 표준을 목록화한다. EU AI Act, 미국 연방/주 규제, 중국 규제, 한국 AI 기본법 및 관련 법령을 포함한다.

 \item \textbf{정보 수집 채널 구축}: 규제 기관 공식 웹사이트, 법률 데이터베이스, 산업 협회, 로펌 뉴스레터, 표준화 기관 발표 등 다양한 채널에서 정보를 수집한다.

 \item \textbf{영향 분석 프로세스}: 규제 변화가 감지되면, 해당 변화가 조직의 AI 시스템에 미치는 영향을 분석한다. 영향의 범위(어떤 AI 시스템이 영향을 받는지), 강도(준수 수준의 변화 정도), 긴급성(시행 시점까지의 잔여 기간)을 평가한다.

 \item \textbf{대응 계획 수립 및 이행}: 영향 분석 결과에 따라 대응 계획을 수립하고, 담당 부서에 할당하며, 이행 상황을 추적한다.

 \item \textbf{정기 보고}: 분기별로 규제 변화 현황과 대응 상황을 경영진에게 보고한다.
\end{enumerate}

%----------------------------------------------------------
\subsection{규제 대응 체크리스트}
%----------------------------------------------------------

\begin{practicaltool}{글로벌 AI 규제 대응 통합 체크리스트}

다국적 기업이 글로벌 AI 규제에 체계적으로 대응하기 위한 통합 체크리스트이다.

\textbf{1. 기반 구축 (Foundation)}
\begin{itemize}
 \item[$\square$] AI 시스템 인벤토리 완성 (전사 AI 시스템 목록화)
 \item[$\square$] 적용 관할권 식별 및 규제 매핑 완료
 \item[$\square$] 글로벌 AI 거버넌스 정책 수립
 \item[$\square$] AI 윤리위원회 또는 거버넌스 위원회 설치
 \item[$\square$] 규제 변화 추적 체계 구축
\end{itemize}

\textbf{2. EU AI Act 대응}
\begin{itemize}
 \item[$\square$] AI 시스템별 리스크 분류(4단계) 완료
 \item[$\square$] 금지 AI 시스템 해당 여부 확인 및 조치
 \item[$\square$] 고위험 AI 시스템 기술 문서 작성
 \item[$\square$] 적합성 평가 절차 준비
 \item[$\square$] EU 대리인(Authorised Representative) 지정
 \item[$\square$] GPAI 모델 해당 시 기술 문서 및 저작권 정책 준비
 \item[$\square$] AI 리터러시 교육 프로그램 수립
\end{itemize}

\textbf{3. 한국 AI 기본법 대응}
\begin{itemize}
 \item[$\square$] 고위험 AI(11개 섹터) 해당 여부 확인
 \item[$\square$] 영향평가 수행 체계 구축
 \item[$\square$] 설명 의무 이행 체계 구축
 \item[$\square$] 생성형 AI 표시 의무 이행 (워터마크 등)
 \item[$\square$] 개인정보보호법 자동화 의사결정 조항 준수
\end{itemize}

\textbf{4. 기타 관할권 대응}
\begin{itemize}
 \item[$\square$] 미국: 주별 규제 동향 추적 및 연방 프리엠션 영향 분석
 \item[$\square$] 중국: 알고리즘 등록 및 콘텐츠 라벨링 의무 확인
 \item[$\square$] 일본: AI 사업자 가이드라인 자발적 준수
 \item[$\square$] 싱가포르: AI Verify 활용 자가 평가
\end{itemize}

\textbf{5. 지속적 관리}
\begin{itemize}
 \item[$\square$] 분기별 규제 격차 분석(Gap Analysis) 수행
 \item[$\square$] 연간 내부 감사 수행
 \item[$\square$] 규제 변화 시 30일 이내 영향 분석 완료
 \item[$\square$] 인시던트 발생 시 72시간 이내 보고 체계 가동
 \item[$\square$] 연간 경영진 보고 및 거버넌스 체계 개선
\end{itemize}
\end{practicaltool}

%==========================================================
\subsection*{제2장 요약}
%==========================================================

본 장에서는 글로벌 AI 거버넌스 프레임워크의 구조와 주요국 규제 동향을 심층 분석하였다. 핵심 내용을 요약하면 다음과 같다.

\begin{enumerate}
 \item \textbf{국제 프레임워크의 계층적 활용}: OECD/UNESCO 원칙을 상위 가이드로, NIST AI RMF를 리스크관리 방법론으로, ISO/IEC 42001을 관리체계 표준으로, IEEE 7000을 설계 표준으로 계층적으로 활용하는 것이 효과적이다.

 \item \textbf{OECD AI 원칙}: 5대 가치 기반 원칙(포용적 성장, 인권 존중, 투명성, 견고성, 책임성)은 글로벌 AI 정책의 기준점이며, 2024년 개정을 통해 생성형 AI와 정보 무결성에 대한 대응이 강화되었다.

 \item \textbf{EU AI Act}: 세계 최초의 포괄적 AI 규제 법률로서, 4단계 리스크 분류(금지/고위험/제한/최소), GPAI 규정, 역외 적용 등을 통해 글로벌 표준을 형성하고 있다. 한국 기업은 EU 시장 진출 시 적합성 평가와 기술 문서 체계 구축이 필수적이다.

 \item \textbf{NIST AI RMF}: 거버넌스(Govern), 매핑(Map), 측정(Measure), 관리(Manage)의 4대 기능과 생성형 AI 프로파일(12대 리스크)을 통해 체계적 리스크관리 방법론을 제공한다.

 \item \textbf{ISO/IEC 42001}: 세계 유일의 인증 가능한 AI 관리체계 표준으로서, 10대 조항 구조와 Annex A 제어 목록을 통해 조직의 AI 거버넌스 체계를 공식화한다. ISO 27001, ISO 9001과의 통합 운영이 효율적이다.

 \item \textbf{한국 AI 기본법}: 11개 섹터의 고위험 AI 분류, 영향평가 의무, 설명 의무, 생성형 AI 표시 의무를 규정하며, 2026년 1월 22일부터 시행된다.

 \item \textbf{규제 대응 전략}: 다국적 기업은 EU AI Act를 글로벌 기본선으로 설정하고, 관할권별 모듈화된 준수 체계를 구축하는 통합 거버넌스 전략이 필요하다. 규제 격차 분석, 규제 매핑 자동화, 규제 변화 추적 체계 등의 실무 도구를 활용하여 체계적으로 대응해야 한다.
\end{enumerate}
